\section{Price extensions and simultaneous minimization}\label{sec:prices}

The compression argument introduces new sites at which to place tensor
legs. A price on the old sites already puts a lower bound on the price
of any proposed placement. Contract the legs at the new sites and
evaluate the tensor that remains. One can do this on any tensor power,
with each new site contracting all its copies at once.
Theorem~\ref{thm:extension} shows that these lower bounds are sharp and
gives a criterion for satisfying several price bounds at the same
placement. In Section~\ref{sec:sync} the tensors to be placed will
depend on a parameter. To control specialization of that parameter,
we prove the criterion over valued fields. For a product of ordinary
wires incident to one new site, the minimum is computed in
\cite[Theorem~\ext{G}{thm:extension}]{oaG}.

\subsection{Integral restrictions}
Let \(K\) be a field with a real additive nonarchimedean valuation
\(\nu\), normalized by \(\nu(0)=+\infty\). Its valuation ring is
\(\mathcal O=\{c\in K:\nu(c)\ge0\}\). We use the given coordinate bases
at all sites. A restriction is \emph{integral} if each local matrix has
entries in \(\mathcal O\). A \emph{valued price} \(p\) of scale \(D>0\) on \(S\)
(here $|S|\ge2$, as before) is real on nonzero tensors, has \(p(0)=-\infty,\ p(\mathbf1)=0\), and
is nonincreasing under integral restriction. Its scalar, product, and
direct-sum rules are
\begin{equation}\label{eq:prices-axioms}
 \begin{gathered}
 p(cT)=p(T)-\nu(c) \quad(c\ne0), \qquad
 p(T\otimes U)=p(T)+p(U),\\
 p\Bigl(\bigoplus_j T_j\Bigr)=D\log\sum_j e^{p(T_j)/D}.
 \end{gathered}
\end{equation}
Here \(\log\) is the natural logarithm. Thus multiplying by a scalar of
positive valuation lowers the price. The trivial valuation recovers
ordinary prices.

We shall construct prices using a convenient equivalent form of the
direct-sum rule. For tensors in common spaces, the right side of
\eqref{eq:prices-axioms} bounds \(p(\sum_j T_j)\) by local addition
of blocks. Conversely, suppose this ordinary sum bound holds and
\(p(\Delta_S(q))=D\log q\) for every \(q\). The block inclusions
express \(\bigoplus_j T_j\) as an ordinary sum and give the desired
upper bound. For the lower bound, take the \(m\)-th power of the direct
sum and project at every site onto the words in which the \(j\)-th
summand occurs \(m r_j\) times. Local permutations identify the
result with \(\bigotimes_jT_j^{\otimes m r_j}\) times a diagonal of
size \(m!/\prod_j(mr_j)!\).
For rational probability vectors \(r\) on the nonzero summands,
take \(m\) through multiples of a common denominator. The price per copy
is at least \(\sum_jr_jp(T_j)-D\sum_j r_j\log r_j\).
Maximization gives the asserted identity.

Write \(q_i\) for the local dimensions of \(T\) and \(\nu_{\min}(T)\)
for the minimum valuation of an entry. For \(T\ne0\),
\begin{equation}\label{eq:prices-range}
 -\nu_{\min}(T)\le p(T)\le -\nu_{\min}(T)+D\sum_i\log q_i.
\end{equation}
For the lower bound project to an entry of minimum valuation. For the
upper bound expand \(T\) in at most \(\prod_iq_i\) coordinate tensors
and use the ordinary sum inequality. Thus prices of fixed
scale form a compact space in the topology of pointwise convergence,
since their axioms are closed in the product of the intervals
\eqref{eq:prices-range}. The wire and transport rules above require only
integral operations and hold also for \(p\). At one site the price of a
nonzero vector is \(-\nu_{\min}\), since after division by a minimum-valuation
entry it is integrally equivalent to the unit. When \(K\) contains a
trivially valued field \(F\), restriction to tensors
over \(F\) gives the ordinary spectral geometry of the wire distances.

\subsection{A test for extending a price}
Fix \(p\) on \(S\), and adjoin a finite set \(L\) of sites. Denote by
\(\widehat p\) an extension to \(S\sqcup L\) of the same scale which
retains \(p\) on tensors with trivial legs at \(L\). Such extensions
exist by placing the added sites at old ones. Restriction to any intermediate
list of sites containing \(S\) gives a valued price of scale \(D\).
Indeed the diagonal on that list restricts to the base diagonal and is a
restriction of the full diagonal, so the diagonal test applies. 

A \emph{closure} applies one integral linear functional at each site of
\(L\) and leaves all base legs free. The \emph{assisted cost} of \(T\) is
\begin{equation}\label{eq:prices-beta}
 \beta_p(T)=\sup_{m\ge1}\frac1m
       \sup_{\ell}p\bigl(\ell(T^{\otimes m})\bigr).
\end{equation}
At a site of dimension \(q\) the functional in this formula acts on
\(K^{q^m}\), so it may couple the \(m\) copies arbitrarily. Every
extension satisfies \(\widehat p(T)\ge\beta_p(T)\) by monotonicity
under closure and additivity on powers.

For example, in the ordinary case take the product of two size-\(q\) wires
from \(a\) to \(y\) and from \(y\) to \(b\), where \(a,b\) are old sites
and \(y\) is new. Already on one copy, closing the two legs at \(y\) by
the pairing \(e_i\otimes e_j\mapsto\delta_{ij}\) recovers the wire from \(a\)
to \(b\). This gives the lower bound \(d(a,b)\log q\) for the product. Placing
\(y\) at \(a\) attains it. More generally, placing \(y\) anywhere on
the geodesic from \(a\) to \(b\) attains the bound by additivity of
distance along that segment.

\begin{theorem}[Price extension]\label{thm:extension}
For every tensor \(T\) on \(S\sqcup L\),
\(\beta_p(T)=\min_{\widehat p}\widehat p(T)\).
For an arbitrary family of nonzero tensors \(T_i\) and real numbers
\(h_i\), a single extension satisfies \(\widehat p(T_i)\le h_i\)
for all \(i\) if and only if
\begin{equation}\label{eq:prices-product-test}
 \beta_p(T_{i_1}\otimes\cdots\otimes T_{i_j})
 \le h_{i_1}+\cdots+h_{i_j}
\end{equation}
for every finite list of indices, with repetitions allowed.
\end{theorem}

\begin{proof}
\emph{Minimizing one tensor.}
The zero tensor is immediate. For a prescribed nonzero \(T\) we first
assign prices to a finite list of distinct nonzero tensors \(A_1,\ldots,A_s\)
containing \(T\). The exponents of their mixed products will serve as
variables. For \(x\in\mathbb Q_{\ge0}^s\) set
\[
 b(x)=\frac1n\beta_p\Bigl(\bigotimes_i A_i^{\otimes nx_i}\Bigr),
 \qquad nx\in\mathbb Z^s.
\]
In \eqref{eq:prices-beta} the supremum can be restricted to multiples
of any given integer by repeating each closure independently. Thus assisted
cost is homogeneous under powers and \(b\) is well defined.
Independent closures, taken at common multiples, make \(b\)
superadditive, hence concave by homogeneity.

The increment per unit of coordinate \(i\) lies in the price interval
\eqref{eq:prices-range} for \(A_i\). For the lower estimate isolate a
minimum-valuation coefficient of every added factor, by coordinate
projections at the base sites as well as the added sites. For the upper
estimate expand each added factor in coordinates. Even for a closure
coupling all copies, fixing the added indices of each term induces an
integral closure on the old factors. The coordinate factors remaining at
base sites have price zero, and the ordinary sum bound pays at most
\(D\sum_j\log q_j\) per added factor of format \((q_j)\).
This proves the increment bounds after taking powers. In particular \(b\)
extends continuously to the real orthant.

At a strictly positive rational \(x\) choose a supporting linear form
\(u\) for this homogeneous concave function, so that
\(b(y)\le u\cdot y\) and \(b(x)=u\cdot x\).
The two-sided increment estimates bound \(u_i\) in the same price
interval. We claim that the numbers \(u_i\) respect the price relations on the
list. For example, if \(A_i=A_j\otimes A_k\), then
\(b(x+t(e_j+e_k-e_i))=b(x)\) for both signs of sufficiently small
rational \(t\). The supporting inequality in these two directions forces
\(u_i=u_j+u_k\). A scalar change shifts \(b\) by minus its valuation
per copy, giving the scalar relation by the same argument. If \(A_i\)
restricts integrally to \(A_j\), replacing copies of \(A_j\) by
copies of \(A_i\) can only increase \(b\). The supporting inequality
in the direction \(e_i-e_j\) therefore gives \(u_i\ge u_j\).
A base tensor increases \(b\) by its exact price per unit,
since its factors pass unchanged through every closure.
For \(\Delta_{S\sqcup L}(q)\), closing its added legs by all-ones
functionals yields \(\Delta_S(q)\), and expanding its \(q\) diagonal
entries gives the reverse bound. Thus its increment is \(D\log q\).

To check the ordinary sum bound, let \(A_i=A_j+A_k\) in the list.
Expand \(n\sigma\) copies of \(A_i\) in a product with multiplicities
\(nx\), for a rational \(0<\sigma\le x_i\). A term choosing \(a\)
copies of \(A_j\) has price after closure at most
\(n u\cdot x-n\sigma u_i+a u_j+(n\sigma-a)u_k\).
The ordinary sum bound on the base and the binomial formula bound
the price after closure of the original tensor by
\[
 n u\cdot x-n\sigma u_i+
 n\sigma D\log\bigl(e^{u_j/D}+e^{u_k/D}\bigr).
\]
Taking the supremum per \(n\) through admissible multiples yields
\(b(x)=u\cdot x\) on the left. This proves the required bound on \(u_i\).
Relations involving zero summands hold automatically.

Let \(x\) tend to the unit vector corresponding to \(T\). The support
vectors have a convergent subsequence and continuity gives
\(u_T=\beta_p(T)\) in the limit. Every finite collection of price
relations and the equality at \(T\) is therefore feasible by applying the
construction to the list of all tensors occurring in those constraints.
Compactness of the product of intervals gives a simultaneous
solution for all such relations. The diagonal test following
\eqref{eq:prices-axioms} makes this solution a valued price. This proves
attainment for one tensor.

\smallskip\noindent
\emph{Satisfying several bounds.}
Any extension obeying the proposed bounds satisfies
\eqref{eq:prices-product-test} by additivity on products. Conversely,
take a finite list of \(k\) bounds. Form \(U\) as a direct sum of
\(N_i\) copies of \(T_i^{\otimes n}\) for each \(i\). In a power of
\(U\) each word of summands is a mixed product. Restricting a closure
to the blocks of that word gives an integral closure, whose price is
bounded by \eqref{eq:prices-product-test}. Sum over the words on the
base and take the supremum per copy to obtain
\[
 \beta_p(U)\le D\log\sum_i N_i e^{nh_i/D}.
\]
Choose \(C>\max_i h_i\) and \(N_i=\lceil e^{n(C-h_i)/D}\rceil\).
These multiplicities bring every \(N_i e^{nh_i/D}\) into the interval
\([e^{nC/D},2e^{nC/D}]\). At an extension minimizing \(U\), the
direct-sum identity then gives
\[
 D\log N_i+n\widehat p(T_i)\le\beta_p(U)\le nC+D\log(2k),
 \qquad
 \widehat p(T_i)\le h_i+\frac{D\log(2k)}{n}.
\]
Let \(n\) tend to infinity and apply compactness, first to the finite
list and then, by the finite intersection property, to the whole family.
\end{proof}

Two consequences will be used below. Given several extensions of \(p\),
there is a single extension whose evaluations are bounded above by any
fixed convex combination of theirs. Indeed, each extension bounds
the assisted cost of every mixed product, so the averaged bounds pass
\eqref{eq:prices-product-test}. We may also enlarge the base before
applying the test. Suppose an enlarged base carries a price agreeing
with \(p\) on the old base. A tensor supported on the old base and the
unknown sites has exactly the same closures, with trivial legs on the
extra base sites. Thus its assisted cost is unchanged. This observation
also applies when the enlarged base has larger scale. Theorem~\ref{thm:extension}
then gives an extension at that larger scale.

\subsection{Geodesic placement}
Assume henceforth in this section that \(K\) contains a trivially valued
subfield \(F\). Given pairs of sites \(a_j,b_j\),
place new sites \(y_j\) first at the \(a_j\) and then at the \(b_j\).
The interpolation consequence of Theorem~\ref{thm:extension} gives an
extension bounding the price of every tensor by the affine interpolation
of these two evaluations, with weights \(1-\theta,\theta\). In particular
\[
 d(a_j,y_j)\le\theta\,d(a_j,b_j),\qquad
 d(y_j,b_j)\le(1-\theta)d(a_j,b_j).
\]
The triangle inequality forces equality in both bounds. Thus \(y_j\)
is the point a fraction \(\theta\) of the way along the geodesic from
\(a_j\) to \(b_j\) in a complete \(\operatorname{CAT}(0)\) spectral
realization over \(F\). The same extension provides the interpolation
bound for all tensors over \(K\).

\begin{lemma}[Geodesic determination and submean inequality]\label{lem:geodesic}
Ordinary price data on any subset of a complete
\(\operatorname{CAT}(0)\) spectral space determine all prices on its
closed geodesic convex hull up to a price-preserving isometry fixing the
given data. Moreover, suppose the prices at a central placement are at
most the weighted average of prices at finitely many end placements, for
every tensor on the original slots. In each placement adjoin points
on corresponding geodesic segments. If the two distances of each new
point to its endpoints at the center are the weighted averages of their
values at the ends, the price inequality holds on the enlarged placements.
\end{lemma}

\begin{proof}
For each new site, prescribe the prices of its size-two wires to its
two endpoints. These prices are the two geodesic distances multiplied
by \(\log 2\). We claim that the price of any \(T\ne0\) on the enlarged
list is the smallest number \(h\) for which the bound \(h\) and all the
wire bounds pass the product test. The actual price passes by closure.
Conversely, enlarge the base to include the actual geodesic points,
and apply Theorem~\ref{thm:extension} to duplicate sites with those
bounds. In a common metric realization, the triangle inequality forces
both endpoint distances of each duplicate to be the prescribed distances.
Uniqueness of geodesics identifies the duplicate with the actual point.
Transport~\eqref{eq:prices-transport}
therefore gives the assertion. For a collapsed base all its prescribed
segments are constant and the assertion follows directly.

Iterating this characterization determines prices, including distances,
on all finite geodesic constructions. Completing their union gives the
closed convex hull by continuity of geodesic interpolation, and
\eqref{eq:prices-transport} extends the price identifications to the
completion. An infinite original subset causes no change since each
finite construction starts from finitely many points.

For the second assertion, a collapsed central placement is immediate.
Otherwise fix a product in the test, a power and a closure on its new
sites. The original submean inequality bounds the central price of the
closed tensor by the average of its closed prices at the ends. At each
end this is at most the price of the unclosed power by restriction,
including a collapsed end, where every nonzero tensor has price zero.
Divide by the power and take the supremum at the center. The averaged
actual upper bounds on the enlarged placements thus pass the product
test at the center, and the averaged wire bounds agree with the
prescribed ones there. Apply the characterization just proved.
\end{proof}

\subsection{Unique evaluations at minimizing sites}
A price minimum can admit several placements, as the two-wire example
shows. To compare minimizers as a parameter varies, we add a strictly
convex penalty at each unknown site. Fix \(o\) on the base. For a new
site \(y\), form the direct sum of \(\mathbf1\) and the size-two wire
from \(o\) to \(y\). Its direct-sum label is visible at every site in
the configuration, so its price is
\[
 f(r)=D\log\bigl(1+\exp(r\log(2)/D)\bigr),\qquad r=d(o,y).
\]
For a nonzero tensor \(U\) on the base and a list of new sites, fix an
integer \(M\ge1\). Its \emph{pin objective} \(V\) is the product of
\(U^{\otimes M}\) with one such penalty tensor for each new site.
For several objectives, each is formed on the base and its own list,
then padded by trivial legs at the other lists.

\begin{lemma}[Pinned evaluations]\label{lem:pinning}
For a finite family of objectives on pairwise disjoint lists of new sites,
there is a same-scale extension minimizing each objective. All such
extensions have the same evaluations on every tensor on the joint list.
These evaluations agree on overlaps between subfamilies. For any such
family, including an infinite one, the resulting evaluations on tensors
over \(F\) are realized in a complete \(\operatorname{CAT}(0)\) spectral
space.

Let \(W\) be the product of a finite family, \(G=\sum_V\beta_p(V)\),
and \(h_T\) the pinned evaluation of a nonzero test tensor \(T\).
For some \(c_D>0\) depending only on \(D\), and every \(n\ge1\),
\begin{equation}\label{eq:prices-pin-limit}
 0\le h_T-\bigl(\beta_p(W^{\otimes n}\otimes T)-nG\bigr)
 \le\frac{\sum_j(\log q_j)^2}{8c_D n},
\end{equation}
where \(q_j\) are the dimensions of the test legs grouped at the added
sites, including dimension one for an unused site.
\end{lemma}

\begin{proof}
Successively enlarge the base and minimize the next objective using
Theorem~\ref{thm:extension}. Each minimization attains its original
minimum by the base-enlargement property.

To compare two assignments within a common extension, interpolate all
corresponding sites at their midpoints, with valued price on \(U\)
bounded by the mean of its two prices. For any two sites \(a,b\) and
their midpoint \(m\),
\begin{equation}\label{eq:prices-midpoint}
 f(d(o,m))\le \tfrac12\bigl(f(d(o,a))+f(d(o,b))\bigr)
                  -c_Dd(a,b)^2 .
\end{equation}
Indeed wire distances are at most \(D\), by restriction of the full
diagonal. Choose \(c_D>0\) small enough that \(f(r)-4c_Dr^2\) is convex
and nondecreasing on \([0,D]\), using the positive lower bounds for
\(f'\) and \(f''\) there. Apply midpoint convexity to this function and
the squared-distance midpoint inequality to \(4c_Dr^2\).
The sites of each objective contain the original base, so its
direct-sum penalty has scale \(D\) even in the larger configuration.
If two assignments minimize that objective, the midpoint is bounded
above by the same minimum minus \(c_D\) times the sum of the squared
distances between corresponding new sites. These distances must all
vanish.

For two minimizing extensions \(P,P'\), enlarge \(P\) by a second copy
of the entire list, imposing upper bounds on its evaluations from \(P'\).
These bounds pass the product test on the old base and therefore on
the enlarged one. Each copied objective attains its minimum, so its sites
are at distance zero from their counterparts. Transport shows
\(P(T)\le P'(T)\) for every test tensor. Interchanging \(P,P'\) gives
equality. Restricting to any subfamily still minimizes its objectives.
The resulting compatible ordinary evaluations on all finite lists
give a spectral space after quotient by zero distances, and hence the
claimed complete extension.

Finally, adjoin an optimizing copy to an arbitrary competing assignment
as above, and let \(d_j\) be the distances between corresponding sites.
The midpoint objective is at least \(G\), and \eqref{eq:prices-midpoint}
gives \(p_{\rm comp}(W)-G\ge 2c_D\sum_jd_j^2\).
Transport gives \(p_{\rm comp}(T)\ge h_T-\sum_jd_j\log q_j\).
Multiply the first inequality by \(n\), add the second, and complete the
squares in the \(d_j\). Taking the minimum by
Theorem~\ref{thm:extension} proves the upper bound in
\eqref{eq:prices-pin-limit}. The optimizing assignment itself gives
the lower bound.
\end{proof}
